%------------------------------------------------------------------------------%
\section{Esboçar o Gráfico da Função Exponencial}

Se quisermos esboçar o gráfico da função $f(x)= a^x$, escolhemos um conjunto de
valores (3 ou 5 pontos) estratégicos para $x$ e montamos uma tabelinha com
esses valores e sua imagem pela função. É comum escolhermos os valores
\[
  x = -2, -1, 0, 1, 2.
\]
Vamos observar os seguintes exemplos.

%------------------------------------------------------------------------------%
\begin{example}
  Esboce o gráfico da função \(f(x)=2^x\).

  \solution
  Para isso, vamos escolher os valores
  \[
    x = -2, -1, 0, 1, 2
  \]
  e montar uma tabela com 3 colunas:
  A primeira coluna esses valores de $x$,
  a segunda coluna com a imagem desses valores pela função $f(x)=2^x$
  e a terceira com o par ordenado $\left(x, f(x)\right)$.
  \begin{ceqn}
  \[
    \begin{array}{r<{\quad}>{\quad}l<{\quad}>{\quad}c}    \hline
       x & f(x) = 2^x                & \left( x, f(x)\right)         \\ \hline
      -2 & f(-2)= 2^{-2}=\frac{1}{4} & \left(-2, \sfrac{1}{4}\right) \\
      -1 & f(-1)= 2^{-1}=\frac{1}{2} & \left(-1, \sfrac{1}{2}\right) \\
       0 & f(0) = 2^0=1              & \left( 0, 1\right)            \\
       1 & f(1) = 2^1=2              & \left( 1, 2\right)            \\
       2 & f(2) = 2^2=4              & \left( 2, 4\right)            \\ \hline
    \end{array}
  \]
  \end{ceqn}
  Inserindo os pontos encontrados no plano cartesiano e sabendo que o
  a forma geral do gráfico da função exponencial obtemos o esboço

  \IncludeGraphicCentredEx{2elevado-1}[0.5]
\end{example}
%------------------------------------------------------------------------------%

%------------------------------------------------------------------------------%
\begin{example}
  Esboce o gráfico da função \(f(x)=\left(\frac{1}{2}\right)^x\).

  \solution
  Para isso, vamos escolher, novamente, os valores
  \[
   x = -2, -1, 0, 1, 2
  \]
  e construir uma tabela com 3 colunas como no exemplo anterior.
  \begin{ceqn}
  \[
    \begin{array}{r<{\quad}>{\quad}l<{\quad}>{\quad}c}    \hline
       x & f( x) = \left(\sfrac{1}{2}\right)^x              & \left( x, f(x)\right)          \\ \hline
      -2 & f(-2) = \left(\sfrac{1}{2}\right)^{-2}=2         & \left(-2, 2\right)             \\
      -1 & f(-1) = \left(\sfrac{1}{2}\right)^{-1}=2         & \left(-1, 1\right)             \\
       0 & f( 0) = \left(\sfrac{1}{2}\right)^0=1            & \left( 0, 1\right)             \\
       1 & f( 1) = \left(\sfrac{1}{2}\right)^1=\sfrac{1}{2} & \left( 1, \sfrac{1}{2}\right)  \\
       2 & f( 2) = \left(\sfrac{1}{2}\right)^2=\sfrac{1}{4} & \left( 2, \sfrac{1}{4}\right)  \\ \hline
    \end{array}
  \]
  \end{ceqn}
  Inserindo os pontos encontrados no plano cartesiano e sabendo que o
  a forma geral do gráfico da função exponencial obtemos o esboço

  \IncludeGraphicCentredEx{meioelevado}[0.5]
\end{example}
%------------------------------------------------------------------------------%

%------------------------------------------------------------------------------%